\documentclass[a4paper,11pt]{report}
\usepackage[utf8]{inputenc}
\usepackage[T1]{fontenc}
\usepackage{lmodern}
\usepackage[ngerman]{babel}
\usepackage[margin=20mm, left=20mm, right=10mm, headheight=15pt, includeheadfoot]{geometry}
\usepackage{fancyhdr}
\usepackage{opensans}
\usepackage{titlesec}
\usepackage{tocloft}
\usepackage{titling}
\usepackage{hyperref}
\usepackage{graphicx}
\usepackage{enumerate}
\usepackage{float}
\usepackage{caption}
\usepackage{listings}
\usepackage{minted}

\usemintedstyle{vs}

% Set default font to OpenSans
\renewcommand*\familydefault{\sfdefault}

% Page style for cover page
\fancypagestyle{cover}{
    \fancyhf{}
    \renewcommand{\headrulewidth}{0pt}
    \renewcommand{\footrulewidth}{0pt}
}

% Page style for main content
\fancypagestyle{main}{
    \fancyhf{}
    \fancyhead[L]{\theauthor}
    \fancyhead[R]{\subject}
    \fancyfoot[C]{\thepage}
    \renewcommand{\headrulewidth}{0.4pt}
    \renewcommand{\footrulewidth}{0pt}
}

\titleformat{\chapter}[block]
  {\normalfont\Large\bfseries} % change \Large to \large or any other size that fits
  {\thechapter}
  {1em}
  {}

  \titlespacing*{\chapter}{0pt}{*4}{*2.5}

% Cover page
\newcommand{\coverpage}{
    \thispagestyle{cover}
    \begin{center}
        {\includegraphics[height=3cm]{fh-logo}}\\[1cm]
        {\LARGE \thetitle}\\[0.5cm]
        {\large \theauthor}\\
        \href{mailto:97hilfel@gmail.com}{97hilfel@gmail.com}\\
    \end{center}
    \tableofcontents
    \clearpage
}

\newcommand{\screenshot}[1]{
    \begin{figure}[H]
        \centering
        \includegraphics[scale=0.375]{#1}
    \end{figure}
}

% Author and subject
\author{Felix Hillebrand}
\newcommand{\subject}{SSD4UE}

\graphicspath{ {./images/} }

% Main document
\begin{document}

\title{Subject~-~Exercise}
\coverpage

\pagenumbering{roman}
\clearpage
\pagenumbering{arabic}
\pagestyle{main}

\chapter{XQuery}
    Übersetzen Sie die nachfolgenden textuellen Angaben in entsprechende XQuery-Ausdrücke,
    sodass die bei jeder Aufgabenstellung angegebene Ausgabe erzeugt wird. 
    Testen Sie diese Ausdrücke auf Basis des bereitgestellten XML-Dokuments \texttt{ue3\_1\_fitnessdokument.xml}. 
    Die Verwendung von XQuery 1.0 ist ausreichend. 
    Der betroffene Namensraum aus dem XML-Dokument soll jeweils beibehalten werden.

    \begin{enumerate}[1.]
    % 1.1
    \item Selektieren Sie die Messwert-Elemente,
    die zu Gewicht-Messungen gehören und sortieren Sie aufsteigend nach dem Datum der Messung.
    Umgeben Sie das gesamte Resultat mit einem Element \texttt{<M>}.
    \textit{(1,5 Punkt)}

    \begin{minted}[frame=lines, breaklines, linenos]{xml}
<M xmlns="http://www.fitness.at/fitness">
    <Messwert typ="Gewicht" wert="80.7" einheit="kg"/>
    <Messwert typ="Gewicht" wert="80.6" einheit="kg"/>
</M>
    \end{minted}
    \textbf{Lösung:}
    \inputminted[frame=lines, breaklines, linenos]{xquery}{assets/ue3_1_1.xq}
    \newpage

    % 1.2
    \item Selektieren Sie Messung-Elemente mit Blutdruck-Werten und verarbeiten Sie nur
    Messungen, die im Kommentar (Notiz) 'maschinell' (case-insensitiv) enthalten. Erstellen
    Sie ein neues Blutdruck-Element mit dem Zeitpunkt der Messung „zeit“ als Attribut.
    Fügen Sie zwei Kind-Elemente hinzu, „Sys“ für den systolischen Wert und „Dia“ für den
    diastolischen Wert. Packen Sie die Ausgabe in ein „RR“-Element. Verwenden Sie zur
    Formatierung der Ausgabe Element- und Attributkonstruktoren. \textit{(1,5 Punkte)}
    \begin{minted}[frame=lines, breaklines, linenos]{xml}
<RR xmlns="http://www.fitness.at/fitness">
    <Blutdruck zeit="2024-01-28T07:55:42">
        <Sys>121.0</Sys>
        <Dia>70.0</Dia>
    </Blutdruck>
    <Blutdruck zeit="2024-01-27T07:48:28">
        <Sys>115.0</Sys>
        <Dia>67.0</Dia>
    </Blutdruck>
</RR>
    \end{minted}
    \textbf{Lösung:}
    \inputminted[frame=lines, breaklines, linenos]{xquery}{assets/ue3_1_2.xq}
    \newpage

    % 1.3
    \item Erstellen Sie eine Aktivitätsübersicht für alle Geräte, mit denen bereits mindestens 2
    Messvorgänge durchgeführt wurden, führen Sie dazu auch die Zeitpunkte an. 
    Erstellen Sie hierfür ein Element '\texttt{Aktivität}' und listen Sie darin '\texttt{Gerät}'-Elemente 
    (mit der '\texttt{id}' als Attribut) auf. 
    Die Kindelemente '\texttt{Online}' von '\texttt{Gerät}' enthalten die Messzeitpunkte. 
    \textit{(1,5 Punkte)}
    \begin{minted}[frame=lines, breaklines, linenos]{xml}
<Aktivität xmlns="http://www.fitness.at/fitness">
    <Gerät id="mg_bd1">
        <Online>2024-01-28T07:55:42</Online>
        <Online>2024-01-27T07:48:28</Online>
    </Gerät>
    <Gerät id="mg_wg1">
        <Online>2024-01-28T08:45:29</Online>
        <Online>2024-01-22T09:05:51</Online>
    </Gerät>
    <Gerät id="mg_bz1">
        <Online>2024-01-28T08:44:47</Online>
        <Online>2024-01-27T07:49:13</Online>
    </Gerät>
</Aktivität>
    \end{minted}
    \textbf{Lösung:}
    \inputminted[frame=lines, breaklines, linenos]{xquery}{assets/ue3_1_3.xq}

    % 1.4
    \item Geben Sie die Titel (wenn vorhanden) und Vor- und Nachnamen der Person in einem Element 'Name' aus. 
    Fügen Sie auch das Geburtsdatum als Element 'Geburtsdatum' hinzu und 
    berechnen Sie auch das Alter (Attribut 'alter'). 
    Umgeben Sie die Elemente mit einem Tag 'Person'. \textit{(2 Punkte)}

    Ergänzen und verwenden Sie für die Berechnung des Alters folgenden Code:

    \begin{minted}[frame=lines, breaklines, linenos]{xquery}
let $dob := (:Date of birth / Geburtsdatum einsetzen:)
let $birth := number(concat(substring($dob, 1, 4), substring($dob, 6, 2),
substring($dob, 9, 2)))
let $now := (:now/ jetzt angeben:)
let $current := number(concat(substring($now, 1, 4), substring($now, 6, 2),
substring($now, 9, 2)))
let $age := floor (($current - $birth) div 10000)
    \end{minted}

    \begin{minted}[frame=lines, breaklines, linenos]{xml}
<Person xmlns="http://www.fitness.at/person">
    <Name>Sandra Sportlich</Name>
    <Geburtsdatum alter="52">1971-08-24</Geburtsdatum>
</Person>
<!-- Titel zum Testen in XML-Dokument einfügen -->
<Person xmlns="http://www.fitness.at/person">
    <Name>Mag. Sandra Sportlich MSc</Name>
    <Geburtsdatum alter="52">1971-08-24</Geburtsdatum>
</Person>
    \end{minted}
    \textbf{Lösung:}
    \inputminted[frame=lines, breaklines, linenos]{xquery}{assets/ue3_1_4.xq}
    \newpage

    % 1.5
    \item Erstellen Sie ein XML-Fragment "Vitalwerte", das alle Messungen mit Messwerten listet.
    Ergänzen Sie die Messwerte um ein Attribut „flag“ wenn der Messwert außerhalb angegebener Grenzwerte liegt. 
    Berücksichtigen Sie nur Grenzwerte für den passenden Mess-Typ und beachten Sie auch das Gültigkeitsdatum. 
    Liegt kein zutreffender Grenzwert vor bzw. ist der gemessene Wert innerhalb der Grenzwerte, entfällt das Attribut "flag".
    Das Attribut enthält "L" für Werte unter dem unteren Grenzwert bzw. "H" für Werte über dem oberen Grenzwert. \textit{(2,5 Punkte)}

    \textit{Hinweis: Verwenden Sie xs:number(...) für die Umwandlung der Grenzwerte zum Wertevergleich.}

    \begin{minted}[frame=lines, breaklines, linenos]{xml}
<Vitalwerte xmlns="http://www.fitness.at/fitness">
    <Messung typ="BlutdruckPuls" zeitpunkt="2024-01-28T07:55:42">
        <Messwert typ="Systolisch" wert="121.0" einheit="mmHg"/>
        <Messwert typ="Diastolisch" wert="70.0" einheit="mmHg"/>
        <Messwert typ="Puls" wert="79.0" einheit="Schläge/min"/>
    </Messung>
    <Messung typ="Gewicht" zeitpunkt="2024-01-28T08:45:29">
        <Messwert typ="Gewicht" wert="80.6" einheit="kg"/>
    </Messung>
        <Messung typ="Schritte" zeitpunkt="2024-01-28T00:00:00">
        <Messwert typ="Schritte" wert="9800" einheit=""/>
    </Messung>
    <Messung typ="BlutdruckPuls" zeitpunkt="2024-01-27T07:48:28">
        <Messwert typ="Systolisch" wert="115.0" einheit="mmHg"/>
        <Messwert typ="Diastolisch" wert="67.0" einheit="mmHg"/>
    </Messung>
    <Messung typ="Gewicht" zeitpunkt="2024-01-22T09:05:51">
        <Messwert typ="Gewicht" wert="80.7" einheit="kg"/>
    </Messung>
        <Messung typ="Schritte" zeitpunkt="2024-01-20T00:00:00">
        <Messwert typ="Schritte" wert="6320" einheit="pro Tag" flag="L"/>
    </Messung>
        <Messung typ="Blutzucker" zeitpunkt="2024-01-28T08:44:47">
        <Messwert typ="Blutzucker" wert="155.5" einheit="mg/dl"/>
    </Messung>
        <Messung typ="Blutzucker" zeitpunkt="2024-01-27T07:49:13">
        <Messwert typ="Blutzucker" wert="201.5" einheit="mg/dl" flag="H"/>
    </Messung>
</Vitalwerte>       
    \end{minted}
    \textbf{Lösung:}
    \inputminted[frame=lines, breaklines, linenos]{xquery}{assets/ue3_1_5.xq}
    \newpage

    \end{enumerate}

    
\chapter{XSLT}
\begin{enumerate}
    % First main item
    \item \hfill % Use hfill to move the figure to the next line
    \begin{description}
        \item[Aufgabe:]
         \hfill
        \item[Lösung:] \hfill \newline % Use hfill to move the figure to the next line
            Solution
    \end{description}
\end{enumerate}
\newpage

\chapter{Legal}
Die Ausarbeitung der Aufgabe wurde durch:
\begin{itemize}
    \item \texttt{OpenAI - GPT-4.5 Turbo}, 
    \item \texttt{OpenAI - GPT-4.5 Vision}, 
    \item \texttt{Anthropic -- Claude 3 Opus}, 
    \item \texttt{Kagi - FastGPT}
\end{itemize}
mit mehreren unterschiedlichen Prompts und Custom Instructions sowie die verwendung des \href{https://xkcd.com/323/}{Ballmer Peak} unterstützt.

\end{document}
